\documentclass[12pt,a4paper]{article}
\usepackage[UTF8]{ctex}
\usepackage{amsmath,amssymb,amsfonts}
\usepackage{graphicx}
\usepackage{geometry}
\usepackage{fancyhdr}
\usepackage{hyperref}
\usepackage{listings}
\usepackage{xcolor}
\usepackage{float}
\usepackage{booktabs}
\usepackage{multirow}
\usepackage{array}

% 页面设置
\geometry{left=2.5cm,right=2.5cm,top=2.5cm,bottom=2.5cm}
\pagestyle{fancy}
\fancyhf{}
\fancyhead[C]{华数杯数学建模竞赛}
\fancyfoot[C]{\thepage}

% 代码样式设置
\lstset{
    basicstyle=\ttfamily\small,
    keywordstyle=\color{blue},
    commentstyle=\color{green!60!black},
    stringstyle=\color{red},
    numbers=left,
    numberstyle=\tiny\color{gray},
    stepnumber=1,
    numbersep=5pt,
    frame=single,
    breaklines=true,
    tabsize=2
}

% 标题信息
\title{华数杯数学建模竞赛论文}
\author{参赛队伍}
\date{\today}

\begin{document}

\maketitle

\newpage
\tableofcontents
\newpage

\section{问题重述}

% 在这里填写问题重述

\section{问题分析}

% 在这里填写问题分析

\section{模型假设}

\begin{enumerate}
    \item 假设1
    \item 假设2
    \item 假设3
\end{enumerate}

\section{符号说明}

\begin{table}[H]
\centering
\begin{tabular}{cl}
\toprule
符号 & 说明 \\
\midrule
$X$ & 变量说明 \\
$Y$ & 变量说明 \\
\bottomrule
\end{tabular}
\caption{符号说明表}
\end{table}

\section{模型的建立与求解}

\subsection{模型一}

% 在这里填写第一个模型

\subsection{模型二}

% 在这里填写第二个模型

\section{模型求解}

% 在这里填写模型求解过程

\section{结果分析}

% 在这里填写结果分析

\section{模型评价与改进}

\subsection{模型优点}

% 在这里填写模型优点

\subsection{模型缺点}

% 在这里填写模型缺点

\subsection{模型改进}

% 在这里填写模型改进方向

\section{参考文献}

\begin{thebibliography}{99}
    \bibitem{ref1} 参考文献1
    \bibitem{ref2} 参考文献2
\end{thebibliography}

\newpage
\appendix
\section{附录}

\subsection{程序代码}

% 在这里可以添加程序代码

\begin{lstlisting}[language=Python]
# 示例代码
def example():
    pass
\end{lstlisting}

\end{document}
